\section*{Bab \ref{ch:b3:relativistic-kinematics} --- Kinematika Relativistik}

\begin{solution}{exo:b3:relativistic-kinematics:1}
$\gamma = 1.005$, $1.155$, $2.294$, $7.09$, $22.4$. Stasiun antariksa: $\beta =
\num{2.57e-5}$, $\gamma - 1 \approx \beta^2/2 = \num{3.3e-10}$; selama
enam bulan ($\num{1.6e7}{}\,\unit{s}$) jam awaknya berjalan
\emph{lambat} sebesar $\approx\qty{5}{ms}$ (hanya efek kecepatannya;
pada ketinggian rendah stasiun itu efek gravitasinya mengurangi tetapi
tidak membalikkan hal ini).
\end{solution}

\begin{solution}{exo:b3:relativistic-kinematics:2}
(a) $\Delta\tau = 2d/c$. (b) Pada tiap separuh detak, fotonnya menempuh
sisi miringnya: $(c\Delta t/2)^2 = (v\Delta t/2)^2 + d^2$ dengan $d =
c\Delta\tau/2$; selesaikanlah: $\Delta t = \Delta\tau/\sqrt{1 - v^2/c^2}$. (c)
Jika $d$ berubah karena geraknya, langkah Pythagorasnya akan keliru. (d)
Biarkanlah dua penggaris serupa berpapasan, masing-masing membawa kuas
di ujungnya yang mengarah ke penggaris satunya. Jika gerak mengerutkan
panjang melintangnya, tiap kerangka akan meramalkan bahwa penggarisnya
sendiri menorehkan tanda \emph{di luar} ujung penggaris satunya --- dua
fakta yang bertentangan tentang goresan kuas yang sama pada \omterm{def:b3:relativistic-kinematics:event}{peristiwa}
berpapasan yang sama. Pertentangan pada satu \omterm{def:b3:relativistic-kinematics:event}{peristiwa} tidak
diperkenankan: panjang melintang tidak dapat berubah.
\end{solution}

\begin{solution}{exo:b3:relativistic-kinematics:3}
$\gamma = 22.4$. (a) $\gamma\tau = \qty{49}{\micro s}$. (b) $v\gamma
\tau = \qty{14.8}{km}$. (c) $15/22.4 = \qty{670}{m}$. (d) Waktu
terbangnya di laboratorium $\qty{50}{\micro s}$: tanpa kerelatifan
$\eu^{-50/2.2} \approx \num{e-10}$ --- tak satu pun; dengan kerelatifan
\omterm{prop:b3:relativistic-kinematics:dilation}{waktu dirinya} $50/22.4 = \qty{2.2}{\micro s}$, dengan bagian $\eu^{-1}
\approx 0.37$.
\end{solution}

\begin{solution}{exo:b3:relativistic-kinematics:4}
(a) $c\Delta t = \qty{600}{m}$, $\Delta x = \qty{300}{m}$: $\Delta s^2
= (600^2 - 300^2)\,\unit{m^2} > 0$, mirip-waktu. (b) Ya: isyarat berlaju
$\Delta x/\Delta t = c/2$ sudah cukup. (c) Laju kerangka yang sama itu, $v =
0.5c$; \omterm{prop:b3:relativistic-kinematics:dilation}{waktu dirinya} $\Delta s/c = \sqrt{600^2 - 300^2}/c = 520/c =
\qty{1.73}{\micro s}$. (d) Kini $c\Delta t = 300 < \Delta x = 600$:
mirip-ruang; tak ada kaitan kausal yang mungkin; tak ada kerangka yang
membawa keduanya ke tempat yang sama, tetapi kerangka pada
$v = c^2\Delta t/\Delta x = 0.5c$ membuat keduanya serentak.
\end{solution}

\begin{solution}{exo:b3:relativistic-kinematics:5}
(a) $\beta = \num{1.29e-5}$: $\gamma - 1 = \num{8.3e-11}$. (b)
$86400\,\unit{s} \times \num{8.3e-11} = \qty{7.2}{\micro s}$ tiap hari.
(c) $c \times \qty{7.2}{\micro s} = \qty{2.2}{km}$ --- navigasinya
akan mati dalam sehari. (d) Neto $+38\,\unit{\micro s}$ berarti
pergeseran biru gravitasinya ($+45.7$) mengalahkan perlambatan
kinematiknya ($-7.2$): pada \qty{20200}{km}, berada lebih tinggi dalam
potensial Bumi mempercepat jamnya lebih daripada laju orbitnya
memperlambatnya.
\end{solution}

\begin{solution}{exo:b3:relativistic-kinematics:6}
(a) $1.6c/1.64 = 0.976c$. (b) $1.8c/1.81 = 0.994c$. (c) $c - u =
\dfrac{(c - u')(c - v)}{c\,(1 + u'v/c^2)}$: kedua faktornya positif,
sehingga $u < c$. (d) Bintiknya adalah \emph{pola} yang bergerak, bukan
sebuah benda: tak ada materi, energi atau informasi yang berjalan dari
satu titik muka Bulan ke titik berikutnya --- tiap fotonnya berangkat ke
Bulan dengan $c$.
\end{solution}

\begin{solution}{exo:b3:relativistic-kinematics:7}
(a) Panjang gelombangnya lebih pendek: mendekat. $f/f_0 = 589/575 = 1.024$;
$(1+\beta)/(1-\beta) = 1.049$: $\beta = 0.024$, sekitar
\qty{7200}{km/s}. (b) Nisbahnya $2$: $(1+\beta)/(1-\beta) = 4$, $\beta =
0.6$. (c) Dengan menjabarkan akarnya, $\Delta\lambda/\lambda \approx
\beta$; pada \qty{600}{nm}, $\qty{1}{\angstrom} = \qty{0.1}{nm}$ memberikan
$\beta = \num{1.7e-4}$: sekitar \qty{50}{km/s} tiap angstrom --- refleks
para astronom. (d) Efek Doppler melintang: \omterm{prop:b3:relativistic-kinematics:dilation}{pemuluran waktu} murni,
berorde $\beta^2$ --- hanya terukur dengan ketelitian atom.
\end{solution}

\begin{solution}{exo:b3:relativistic-kinematics:8}
(a) $20/2 = \qty{10}{m}$: ia muat, tepat pas, dan kedua pintunya dapat
ditutup serentak (kerangka lumbung). (b) Dalam kerangka pelarinya,
\omterm{def:b3:relativistic-kinematics:event}{peristiwa} ``pintu depan menutup'' dan ``pintu belakang membuka''
\emph{tidak} serentak: pintu keluarnya membuka sebelum pintu masuknya
menutup, dan galah \qty{20}{m} itu menyusup melalui lumbung \qty{5}{m}
tanpa pernah terkurung. (c) Kerangka lumbung: $\Delta t = 0$ pada $\Delta x =
\qty{10}{m}$. Kerangka pelari: $|\Delta t'| = \gamma v\Delta x/c^2 =
2 \times 0.866c \times 10/c^2 = \qty{58}{ns}$. (d) ``Galahnya
terkurung'' diam-diam menegaskan bahwa dua \omterm{def:b3:relativistic-kinematics:event}{peristiwa} terpisah (kedua
pintu tertutup) serentak --- pernyataan yang bergantung kerangka, bukan
sebuah fakta.
\end{solution}

\begin{solution}{exo:b3:relativistic-kinematics:9}
$\gamma = 1.25$ pada $0.6c$. (a) Panjang bergeraknya $100/1.25 = \qty{80}{m}$
pada \qty{1.8e8}{m/s}: \qty{444}{ns}. (b) $240/\num{1.8e8} =
\qty{1.33}{\micro s}$. (c) $|\Delta t'| = \gamma v\Delta x/c^2 =
\qty{200}{ns}$; dari $t' = \gamma(t - vx/c^2)$, penjepit pada $x$ yang
lebih besar --- yakni yang di depan --- menembak \emph{lebih dahulu}
bagi roketnya. (d)
Stasiun: $\Delta s^2 = 0 - 80^2 = -6400\,\unit{m^2}$. Roket: $\Delta
x' = \gamma\,\Delta x = \qty{100}{m}$, $c\Delta t' = \qty{60}{m}$:
$60^2 - 100^2 = -6400\,\unit{m^2}$ --- invarian.
\end{solution}

\begin{solution}{exo:b3:relativistic-kinematics:10}
(a) Terra: \qty{20}{yr}; Stella: $20/\gamma = \qty{12}{yr}$. (b)
Jaraknya terkerut menjadi $8/\gamma = \qty{4.8}{ly}$, yang ditempuh dalam $4.8/0.8
= \qty{6}{yr}$ waktunya --- selaras dengan \qty{12}{yr} untuk
perjalanan pulang pergi. (c) \omterm{def:b3:relativistic-kinematics:event}{Keserentakan} dengan \omterm{def:b3:relativistic-kinematics:event}{peristiwa} pembalikan $(t =
\qty{10}{yr},\ x = \qty{8}{ly})$: kerangka keberangkatannya memberi jam
Bumi $t - vx/c^2 = 10 - 6.4 = \qty{3.6}{yr}$; kerangka kepulangannya, $10 +
6.4 = \qty{16.4}{yr}$. ``Sekarang di Bumi'' menurutnya melompat sebesar \qty{12.8}{yr}
pada tekukan itu; pembukuannya $3.6 + 12.8 + 3.6 = \qty{20}{yr}$ cocok
persis dengan Terra. (d) Stella menempati \emph{dua} \omterm{thm:b3:relativistic-kinematics:postulates}{kerangka inersial}
yang disambung oleh percepatan yang dapat dirasakannya; Terra menempati
satu. Garis dunia lurus antara \omterm{def:b3:relativistic-kinematics:event}{peristiwa} keberangkatan dan pertemuan
kembali membawa \omterm{prop:b3:relativistic-kinematics:dilation}{waktu diri} terpanjang; hanya garis Stella yang tertekuk.
\end{solution}

\begin{solution}{exo:b3:relativistic-kinematics:11}
(a) Jika pemetaannya taklinear, gerak seragam tidak akan tampak seragam
dalam kerangka yang lain, sehingga titik-titik ruang dan waktu yang
homogen menjadi terbedakan. (b) Dengan menyulihkan satu hubungan ke
hubungan yang lain: $t' = \gamma t +
(1 - \gamma^2)x/\gamma v$. (c) Dengan mengambil $x = ct$, $x' = ct'$:
$\gamma(c - v)t = c\gamma t + c(1 - \gamma^2)ct/\gamma v$, yang
terselesaikan menjadi $\gamma^2 = 1/(1 - v^2/c^2)$. (d) Asas
kerelatifan memberikan $\gamma$ yang sama untuk kedua arahnya (tak ada
kerangka istimewa); postulat cahayanya mengubah fungsi bebas yang
tersisa menjadi $\gamma(v)$ yang tertentu.
\end{solution}

\begin{solution}{exo:b3:relativistic-kinematics:12}
(a) Dengan $\cosh\varphi = \gamma$, $\sinh\varphi = \gamma\beta$,
transformasinya persis rotasi hiperbolik yang dinyatakan itu, dan
$\cosh^2 - \sinh^2 = 1$ mempertahankan $c^2t^2 - x^2$. (b)
$\tanh(\varphi_1 + \varphi_2) = (\tanh\varphi_1 + \tanh\varphi_2)/(1 +
\tanh\varphi_1\tanh\varphi_2)$: tepat hukum penggabungannya ---
kecepatan tidak dijumlahkan, rapiditas dijumlahkan. (c) $\varphi = g\tau/c$, sehingga
$\beta = \tanh(g\tau/c)$; $\operatorname{artanh}(0.99) = 2.65$ memberikan
$\tau = 2.65c/g \approx \qty{8.1e7}{s} \approx 2.6$ tahun waktu
kapalnya. (d) $\varphi$ menjangkau seluruh bilangan real sedangkan
$\tanh\varphi$ jenuh pada $1$: kita dapat dipercepat selamanya, dengan
rapiditas bertambah secara linear, sementara lajunya hanya merayap
menuju $c$.
\end{solution}

\begin{solution}{pb:b3:relativistic-kinematics:1}
\textbf{1.} Semua muon benar-benar serupa dan meluruh dengan laju
statistik yang tetap: bagian populasi yang bertahan mengukur \omterm{prop:b3:relativistic-kinematics:dilation}{waktu diri}
yang berlalu sepasti jarum jam.
\textbf{2.} $c\tau = \qty{3.00e8}{} \times \qty{2.2e-6}{} =
\qty{660}{m}$.
\textbf{3.} $t = \qty{15}{km}/c = \qty{50}{\micro s}$: bagiannya
$\eu^{-50/2.2} = \eu^{-22.7} \approx \num{1.4e-10}$.
\textbf{4.} Partikel yang ``tidak dapat'' menempuh lebih dari satu
kilometer justru melintasi lima belas kilometer dan tiba dalam jumlah
besar.
\textbf{5.} Telapak tangan $\sim\qty{100}{cm^2}$: beberapa muon tiap
sekon; tubuh dengan penampang mendatar $\sim\num{e3}\,\unit{cm^2}$:
berorde lima belas di antara dua denyut jantung --- kerelatifan turun
bagai hujan pada setiap orang, selalu.
\textbf{6.} $\gamma = 1/\sqrt{1 - 0.995^2} = 10.0$.
\textbf{7.} Membandingkan dua cacah dari populasi terpilih yang
\emph{sama} pada beda ketinggian yang diketahui menyingkirkan setiap
anggapan tentang di mana dan berapa banyak muon dilahirkan.
\textbf{8.} $t = 1907/(0.995 \times \num{3e8}) = \qty{6.39}{\micro s}$.
\textbf{9.} $\eu^{-6.39/2.2} = 0.055$: sekitar $31$ tiap jam.
\textbf{10.} $t/\gamma = \qty{0.64}{\micro s}$.
\textbf{11.} $\eu^{-0.64/2.2} = 0.75$: sekitar $420$ tiap jam.
\textbf{12.} Yang terukur $408 \pm 9$: nilai kerelatifan $\approx 420$
sesuai, mengingat muonnya sedikit melambat dalam penyerap detektornya;
nilai fisika klasik $31$ meleset sebesar faktor tiga belas. Gunungnya
yang memutuskan.
\textbf{13.} $408/563 = 0.725 = \eu^{-t_{\text{diri}}/\tau}$:
$t_{\text{diri}} = \qty{0.71}{\micro s}$, sehingga $\gamma_{\exp} =
6.39/0.71 = 9.0$ --- tepat di dalam $8.8 \pm 0.8$ milik Frisch dan Smith.
\textbf{14.} $1907/10 = \qty{191}{m}$.
\textbf{15.} $191/(0.995c) = \qty{0.64}{\micro s}$: peluruhan $25\%$
yang sama --- yakni kisah muon tentang cacah yang sama.
\textbf{16.} Cacah detak adalah himpunan \omterm{def:b3:relativistic-kinematics:event}{peristiwa} berimpit setempat;
\omterm{def:b3:relativistic-kinematics:event}{peristiwa} beserta cacahnya bersifat invarian terhadap kerangka ---
kerangka boleh berselisih tentang waktu dan panjang, tak pernah tentang
apa yang terbaca pada sebuah pencacah.
\textbf{17.} $\Delta s^2 = (c \times \qty{6.39}{\micro s})^2 -
(\qty{1907}{m})^2 = (1917^2 - 1907^2)\,\unit{m^2} = (\qty{196}{m})^2$:
$\Delta s/c = \qty{0.65}{\micro s}$ --- yakni \omterm{prop:b3:relativistic-kinematics:dilation}{waktu dirinya}, yang
dihitung tanpa pernah meninggalkan kerangka gunungnya.
\textbf{18.} Laju peluruhannya akan bergantung pada ketinggian untuk
setiap laju secara sama; kenyataannya kebertahanannya mengikuti $\gamma$
--- pilihan yang lebih lambat memulur lebih sedikit, tepat seperti yang
ditetapkan $\gamma(v)$.
\textbf{19.} Umur sebesar $\gamma\tau$ sampai satu per seribu (cincin
penyimpan muon di CERN); cincinnya menjadikan muon seorang pengelana
yang terus-menerus dipercepat --- ``saudara kembar yang bepergian''
tetap muda bahkan ketika perjalanannya adalah pembalikan tanpa akhir.
\textbf{20.} $\beta = \num{8.3e-7}$: $\Delta t = \tfrac12\beta^2
\times \qty{144000}{s} = \qty{50}{ns}$.
\textbf{21.} Laju pesawatnya menambah atau mengurangi laju Bumi yang
berputar, sedangkan efek ketinggian (gravitasi) sama untuk kedua arah:
penerbangan ke timur dan ke barat memisahkan kedua sumbangan itu secara
bersih (hasil 1971 cocok dengan keduanya).
\textbf{22.} Dari \cref{exo:b3:relativistic-kinematics:5}:
\qty{7.2}{\micro s} tiap hari.
\textbf{23.} $7 \times 7.2\,\unit{\micro s} \times c \approx
\qty{15}{km}.$
\textbf{24.} Posisinya akan menghanyut ratusan meter dalam hitungan
jam, dan kilometer dalam sehari --- navigasinya akan rusak secara kasat
mata pada pagi pertama.
\textbf{25.} Pada 1963 sebuah gunung, dua pencacah dan jam
\qty{2.2}{\micro s} mengukur $\gamma_{\exp} = 9.0$ terhadap ramalan $10$
(dengan perlambatannya, $8.8 \pm 0.8$): \omterm{prop:b3:relativistic-kinematics:dilation}{pemuluran waktu} terbenarkan
dalam selisih $10\%$ --- dan fisika yang sama dikoreksi diam-diam, tiap
sekon, dalam tiap satelit navigasi di atas kepala Anda.
\end{solution}
